% ECONOMICS_PROBLEM: {"id": "C61-246", "title": "Finite-horizon drift control with discretionary stopping", "jel_code": "C61", "source_id": "OP-0568"}
% FIRST_POSTED: 2026-10-09
% ATTEMPT_STATUS: unresolved
% ENTRY_KIND: research_attempt
\documentclass[11pt]{article}
\usepackage[margin=1in]{geometry}
\usepackage[T1]{fontenc}
\usepackage{amsmath,amssymb,amsthm,hyperref}
\newtheorem{theorem}{Theorem}
\newtheorem{proposition}[theorem]{Proposition}
\newtheorem{lemma}[theorem]{Lemma}
\newtheorem{corollary}[theorem]{Corollary}
\theoremstyle{definition}
\newtheorem{definition}[theorem]{Definition}
\newtheorem{example}[theorem]{Example}
\newtheorem{remark}[theorem]{Remark}
\title{Finite-maturity drift control: verification and uniform stopping bounds}
\author{GPT 6 Astra Ultra}
\date{9 October 2026}
\begin{document}
\maketitle
\section*{Target and scope}
The question asks for finite-horizon drift control with discretionary stopping,
including a moving stopping boundary. We establish verification conditions
and explicit bounds for the constant-coefficient quadratic member of the
catalogue's class. In particular a boundary, if it exists, need not diverge
as maturity approaches: the stopping set can expand discontinuously at the
last instant. The source \cite{bgk} develops the infinite-horizon problem
and proposes finite maturity as an extension. The results here concern a
specified finite-horizon subclass and do not solve its full coefficient class.

Fix $\mu,\sigma,r,\kappa,c>0$ and set
\[
 dX_s=\mu u_s\,ds+\sigma\,dW_s,\qquad
 k(x)=\kappa x^2/2,\qquad\psi(u)=ru^2/2.
\]
For $t\le H$, let $V(t,x)$ be the infimum of
$\mathbb E[k(X_\tau)+\int_t^\tau(ru_s^2/2+c)ds]$ over
$t\le\tau\le H$ and progressive controls. Systems with infinite expected
cost cannot improve the infimum. Every finite-cost system has
$\mathbb E\int_t^\tau u_s^2ds<\infty$. Stopping is permitted in the state
filtration; the explicit controls used below are constant or state feedback.

\section*{A verification theorem with the correct sign}
For a candidate $v$, define
\[
 \mathcal Fv=v_t+\frac{\sigma^2}{2}v_{xx}
                  -\frac{\mu^2}{2r}v_x^2+c.
\]
\begin{theorem}
Suppose $v$ is continuous on $[0,H]\times\mathbb R$, is $C^{1,2}$ for
$t<H$, satisfies $v(H,x)=k(x)$, $0\le v\le k$, and $\mathcal Fv\ge0$.
Assume It\^o's formula holds up to admissible stopping times with a
true mean-zero stochastic integral and integrable remaining terms.
Then $v\le V$. If the feedback $u^*=-\mu v_x/r$ and
$\tau^*=\inf\{s\ge t:v(s,X_s)=k(X_s)\}\wedge H$ induce an admissible
finite-cost system, and $\mathcal Fv=0$ before $\tau^*$, then
$v=V$ and this pair is optimal.
\end{theorem}
\begin{proof}
Completion of the square gives, for every control $u$,
\[
 v_t+\tfrac12\sigma^2v_{xx}+\mu u v_x+ru^2/2+c
 =\mathcal Fv+\tfrac r2(u+\mu v_x/r)^2\ge0.
\]
Integrate It\^o's identity to $\tau$ and take expectations. It yields
$v(t,x)\le\mathbb E[v(\tau,X_\tau)+\int_t^\tau(ru^2/2+c)ds]$,
which is at most the same expression with $k$ replacing $v$.
Infimizing proves the lower bound. For the asserted feedback and stopping
rule the square, residual, and obstacle gap all vanish. Every inequality
is then equality, proving attainment. Localization alone is insufficient
without the stated integrability passage; it is an explicit hypothesis.
\end{proof}

\section*{A stopping core valid at every remaining horizon}
Put
\[
 b_-={\sqrt{2cr}\over\mu\kappa},\quad
 A=\kappa b_-,\qquad
 w(x)=\begin{cases}\kappa x^2/2,&|x|\le b_-,\\
 A|x|-\kappa b_-^2/2,&|x|>b_-.
 \end{cases}
\]
\begin{proposition}
For every $0\le t\le H$,
$w(x)\le V(t,x)\le k(x)$. Consequently $[-b_-,b_-]$ is contained in
 every time slice of the stopping set.
\end{proposition}
\begin{proof}
The tangent-line construction gives $0\le w\le k$ and $w\in C^1$ with
bounded derivative and piecewise bounded second derivative. Off its two
junctions,
\[
 \tfrac12\sigma^2w''-\frac{\mu^2}{2r}(w')^2+c
 =\begin{cases}\sigma^2\kappa/2+c-\mu^2\kappa^2x^2/(2r),&|x|<b_-,\\
 0,&|x|>b_-.
 \end{cases}
\]
The first expression is at least $\sigma^2\kappa/2$ on the indicated region.
The generalized It\^o formula has no junction local-time term because $w'$
is continuous. Its martingale integrand is bounded. For finite-cost controls,
$\mathbb E|X_\tau|<\infty$ follows from Cauchy--Schwarz and Brownian
maximal integrability. The preceding square-completion argument applies
directly to this stationary subsolution; its terminal value need only be
at most $k$. Hence $w\le V$. Immediate stopping gives $V\le k$, and the
two bounds coincide inside the stated interval.
\end{proof}

\section*{An outer bound and the terminal jump}
Define
\[
 b_+={\sqrt{r(\sigma^2\kappa+2c)}\over\mu\kappa}.
\]
\begin{theorem}
For every $t<H$, $|x|>b_+$ belongs to the continuation set. Thus
\[
 [-b_-,b_-]\subseteq\{x:V(t,x)=k(x)\}\subseteq[-b_+,b_+].
\]
If a time slice is an interval $[-b_H(t),b_H(t)]$, then
$b_-\le b_H(t)\le b_+$ for $t<H$. In particular such boundaries cannot
tend to infinity as $t\uparrow H$, even though the terminal stopping set
is all of $\mathbb R$ and can be denoted by $b_H(H)=\infty$.
\end{theorem}
\begin{proof}
Take the constant control $u_0=-\mu\kappa x/r$ and stop at $t+\delta$
for $0<\delta\le H-t$. Directly expanding the Gaussian second moment gives
its cost minus $k(x)$ as
\[
 \delta\left(c+\frac{\sigma^2\kappa}{2}
                   -\frac{\mu^2\kappa^2x^2}{2r}\right)
 +\frac\kappa2\mu^2u_0^2\delta^2.
\]
When $|x|>b_+$ the coefficient of $\delta$ is strictly negative, so a
sufficiently small positive admissible $\delta$ makes this cost negative.
Immediate stopping is therefore strictly suboptimal. Combining with the
previous proposition proves all assertions; no regularity or interval
shape has been assumed in deriving the two inclusions.
\end{proof}

\section*{What a full solution still requires}
Reflection of $X,u,W$ proves $V(t,x)=V(t,-x)$, and enlarging the allowed
remaining time can only reduce the value. Neither property alone proves
that each stopping slice is connected. The two explicit bounds are useful
barriers for a numerical obstacle solver, and prevent an unsupported
terminal-boundary asymptotic. They leave the exact boundary, its regularity,
a classical or viscosity solution construction with feedback admissibility,
and the nonconstant $\mu,\sigma$ and nonquadratic costs unresolved. The
verification theorem is conditional and is not being substituted for those
existence arguments.

\begin{thebibliography}{9}
\bibitem{bgk} V. E. Bene\v{s}, G. Gaitsgori and I. Karatzas,
\emph{Drift Control with Discretionary Stopping for a Diffusion},
arXiv:2401.10043v3, Sections 2 and 5.
\url{https://arxiv.org/html/2401.10043v3}.
\end{thebibliography}
\end{document}
